% Options for packages loaded elsewhere
\PassOptionsToPackage{unicode}{hyperref}
\PassOptionsToPackage{hyphens}{url}
\PassOptionsToPackage{dvipsnames,svgnames,x11names}{xcolor}
%
\documentclass[
  10pt,
]{article}
\usepackage{amsmath,amssymb}
\usepackage{iftex}
\ifPDFTeX
  \usepackage[T1]{fontenc}
  \usepackage[utf8]{inputenc}
  \usepackage{textcomp} % provide euro and other symbols
\else % if luatex or xetex
  \usepackage{unicode-math} % this also loads fontspec
  \defaultfontfeatures{Scale=MatchLowercase}
  \defaultfontfeatures[\rmfamily]{Ligatures=TeX,Scale=1}
\fi
\ifPDFTeX\else
  % xetex/luatex font selection
\fi
% Use upquote if available, for straight quotes in verbatim environments
\IfFileExists{upquote.sty}{\usepackage{upquote}}{}
\IfFileExists{microtype.sty}{% use microtype if available
  \usepackage[]{microtype}
  \UseMicrotypeSet[protrusion]{basicmath} % disable protrusion for tt fonts
}{}
\makeatletter
\@ifundefined{KOMAClassName}{% if non-KOMA class
  \IfFileExists{parskip.sty}{%
    \usepackage{parskip}
  }{% else
    \setlength{\parindent}{0pt}
    \setlength{\parskip}{6pt plus 2pt minus 1pt}}
}{% if KOMA class
  \KOMAoptions{parskip=half}}
\makeatother
\usepackage{xcolor}
\usepackage[margin=0.8in]{geometry}
\usepackage{longtable,booktabs,array}
\usepackage{calc} % for calculating minipage widths
% Correct order of tables after \paragraph or \subparagraph
\usepackage{etoolbox}
\makeatletter
\patchcmd\longtable{\par}{\if@noskipsec\mbox{}\fi\par}{}{}
\makeatother
% Allow footnotes in longtable head/foot
\IfFileExists{footnotehyper.sty}{\usepackage{footnotehyper}}{\usepackage{footnote}}
\makesavenoteenv{longtable}
\setlength{\emergencystretch}{3em} % prevent overfull lines
\providecommand{\tightlist}{%
  \setlength{\itemsep}{0pt}\setlength{\parskip}{0pt}}
\setcounter{secnumdepth}{-\maxdimen} % remove section numbering

\usepackage{fontspec}
\setmainfont{DejaVu Serif}[Scale=0.92]
\setsansfont{DejaVu Sans}[Scale=0.92]
\setmonofont{DejaVu Sans Mono}[Scale=0.84]
\usepackage{etoolbox}
\AtBeginEnvironment{longtable}{\scriptsize\setlength{\tabcolsep}{3pt}}
\usepackage{microtype}
\usepackage{newunicodechar}
\newunicodechar{⋊}{\ensuremath{\rtimes}}
\newunicodechar{∎}{\ensuremath{\square}}
\newunicodechar{ℓ}{\ensuremath{\ell}}
\setlength{\emergencystretch}{3em}
\renewcommand{\arraystretch}{1.12}
\ifLuaTeX
  \usepackage{selnolig}  % disable illegal ligatures
\fi
\IfFileExists{bookmark.sty}{\usepackage{bookmark}}{\usepackage{hyperref}}
\IfFileExists{xurl.sty}{\usepackage{xurl}}{} % add URL line breaks if available
\urlstyle{same}
\hypersetup{
  pdftitle={Beyond the 2013 table: first values for 64 cells of the bipartite degree--diameter problem},
  pdfauthor={Doğa Özden, with Claude Code (lab notes, job paper-beyond)},
  colorlinks=true,
  linkcolor={Maroon},
  filecolor={Maroon},
  citecolor={Blue},
  urlcolor={Blue},
  pdfcreator={LaTeX via pandoc}}

\title{Beyond the 2013 table: first values for 64 cells of the bipartite degree--diameter problem}
\author{Doğa Özden, with Claude Code (lab notes, job paper-beyond)}
\date{November 1, 2026 (draft)}

\begin{document}
\maketitle

\hypertarget{what-this-note-covers}{%
\section{1. What this note covers}\label{what-this-note-covers}}

N\textsuperscript{b}(Δ, D) is the largest number of vertices of a bipartite graph with maximum degree Δ and diameter
at most D. The 2013 survey table (Miller--Širáň, DS14v2, as reproduced on the Combinatorics Wiki) lists lower bounds
for Δ ≤ 16 and D ≤ 10. Between October 26 and November 1, 2026, cloud jobs of the lab built and checked graphs for
cells outside that table. This note gives one value for each of 64 of them:

\begin{itemize}
\tightlist
\item
  Δ = 17, 18, 19, 20 at D = 5, \ldots, 10 (24 cells);
\item
  D = 11 at Δ = 3, \ldots, 16 and D = 12 at Δ = 3, \ldots, 14 (26 cells);
\item
  D = 13 at Δ = 3, \ldots, 10 and D = 14 at Δ = 3, \ldots, 8 (14 cells).
\end{itemize}

The Δ = 17..20 cells at D = 3 and 4 were filled in the same week but are not in the table below: at Δ = 17, 18, 20
they are incidence graphs of PG(2, Δ−1) and W(Δ−1) (Moore graphs, so optimal), and at Δ = 19 they are the lab's
fourth-hat graph (624) and the AS(17) graph of Corollary A1 (10,982) (both in
\texttt{papers/\allowbreak{}Ozden\_\allowbreak{}polished\_\allowbreak{}results.pdf}), checked all-pairs by the coordinator on October 26.

Words used below: a \emph{hole} is a group element (or vertex pair) not yet within the required distance, and the
\emph{cost} of a candidate is its number of holes; \emph{far pairs} are vertex pairs at distance more than D; an
\emph{evaluation} is one diameter test inside a search; \emph{squad cards} (r2/e16 and so on) are short runs by Claude
Haiku helpers, each trying one idea, directed by an Opus session (\texttt{research/\allowbreak{}squads/\allowbreak{}LOG.md}).

Each value is a lower bound: one explicit graph, built from files in the repository. All but three were checked
exactly by a program that does not share code with the search that found them; the three level graphs at D = 7 rest on
a theorem (section 4). A literature card of the new-cells-beyond job (\texttt{research/\allowbreak{}squads/\allowbreak{}beyond/\allowbreak{}r2/\allowbreak{}d10/\allowbreak{}LIT.md} on
its branch) looked for published bipartite values in 40 of these 64 cells (Δ = 17..20 at D = 5..10, and D = 11, 12 at
Δ = 3..10; its list of 48 also has the D = 3, 4 cells) and found none, but it could not read the Combinatorics Wiki
page (HTTP 503) or the tables of some papers. For the other 24 cells no literature search was made (apart from the
cubic censuses below). So we call the values \emph{first values for the lab}, not records.

\textbf{Cubic cells: published graphs.} The census of cubic edge-transitive graphs (fostercensus.graphsym.net; CAT =
arc-transitive, CSS = semisymmetric) lists bipartite cubic graphs of diameter 11 on 1,200 vertices (CAT(1200,1)), 12
on 2,304 (CSS(2304,5)), 13 on 3,360 (CAT(3360,4)) and 14 on 6,272 (CSS(6272,7)), and the census of cubic
vertex-transitive graphs one of diameter 11 on 1,250 (lab notes:
\texttt{research/\allowbreak{}squads/\allowbreak{}own-small/\allowbreak{}newcells/\allowbreak{}README.md}, where the 1,200-, 1,250- and 2,304-vertex graphs were also
checked all-pairs). All are below the lab's values at D = 11, 12 (1,474 and 2,368). But at (3,14) a published graph,
CSS(6272,7), is larger than the lab's 5,706, and at (3,13) the lab's 3,400 is only 40 above CAT(3360,4). The (3,13)
and (3,14) census graphs have not been checked here, and censuses of larger orders were not searched.

For comparison, the bipartite Moore bound is M\textsuperscript{b}(Δ, D) = 2(1 + (Δ−1) + \ldots{} +
(Δ−1)\textsuperscript{D−1}), and for Δ ≥ 3, 5 ≤ D ≤ 187, D ≠ 6, N\textsuperscript{b}(Δ, D) ≤ M\textsuperscript{b}(Δ,
D) − 6 (Feria-Purón and Pineda-Villavicencio, Discrete Appl. Math. 160 (2012); Prop. 6.4 in arXiv 1010.5651v2, the
version that was read). Our values are 100 \% of the Moore bound at (17,6), (18,6), (20,6), 80 \% at (19,6), 24--27 \%
at D = 5, 14--16 \% at D = 7, 4.5--7.7 \% for the orbit-S Haar graphs at D = 8..14 (0.8--3.1 \% in the four cells
stopped at the jobs' size cap), and 7--36 \% for the metacyclic Haar graphs and lifts at Δ = 3..6.

\hypertarget{where-the-numbers-come-from}{%
\section{2. Where the numbers come from}\label{where-the-numbers-come-from}}

Every value in the table was copied from one of these, and the vertex counts were re-derived from the group orders or
formulas (\texttt{python3\ research/\allowbreak{}paper/\allowbreak{}make\_\allowbreak{}beyond.py} prints the comparison and rewrites the table in this file):

\begin{longtable}[]{@{}
  >{\raggedright\arraybackslash}p{(\columnwidth - 4\tabcolsep) * \real{0.2087}}
  >{\raggedright\arraybackslash}p{(\columnwidth - 4\tabcolsep) * \real{0.3391}}
  >{\raggedright\arraybackslash}p{(\columnwidth - 4\tabcolsep) * \real{0.4522}}@{}}
\toprule\noalign{}
\begin{minipage}[b]{\linewidth}\raggedright
job (date)
\end{minipage} & \begin{minipage}[b]{\linewidth}\raggedright
branch
\end{minipage} & \begin{minipage}[b]{\linewidth}\raggedright
cells
\end{minipage} \\
\midrule\noalign{}
\endhead
\bottomrule\noalign{}
\endlastfoot
new-cells-beyond (Oct 26) & claude/new-cells-beyond-etos3t & Δ = 17..20 at D = 3..10; first values at D = 11, 12 for Δ
= 3..10 \\
conics-2 (Oct 29--30) & claude/conics-2-diameter-7-ic0oat & (17,7) \\
own-small (Oct 26) & claude/own-small-cells-dkwo42 & (3,12) \\
beyond-d11 (Oct 31--Nov 1) & claude/new-cells-d11-d12-beyond-8sas9r & D = 11 at Δ = 11..14, D = 12 at Δ = 11..13;
raises at Δ = 4..10 \\
beyond-d11b (Nov 1) & claude/beyond-d11-d12-new-cells-3t3djx & (15,11), (16,11), (14,12); raises of (12,12), (13,12)
and of Δ = 4..6 \\
beyond-d13 (Nov 1) & claude/beyond-d13-search-z35z61 & D = 13 at Δ = 3..10, D = 14 at Δ = 3..8; raise of (4,12) \\
beyond-small (Nov 1) & claude/beyond-small-cells-fwb6si & (17,5); exact checks at D = 6 \\
\end{longtable}

(Dates are when the jobs ran; the harvest and the coordinator recorded the values on master between October 26 and
November 1.) On master the values live in \texttt{research/\allowbreak{}RESULTS.md} (harvest and coordinator lines), the harvest
check reports \texttt{research/\allowbreak{}frontier/\allowbreak{}runs/\allowbreak{}harvest\_\allowbreak{}check\_\allowbreak{}2026-10-31.txt} and \texttt{…\_\allowbreak{}2026-11-01.txt}, and the
coordinator's logs \texttt{research/\allowbreak{}coordinator\_\allowbreak{}check/\allowbreak{}check\_\allowbreak{}2026-10-29.txt} \ldots{} \texttt{check\_\allowbreak{}2026-11-01.txt}.
RESULTS.md has no line yet for (17,5) 35,618, (3,13), (3,14), (4,12) 74,066, (4,13) or (4,14); on master these are in
the coordinator logs of October 31 and November 1, with their edge lists. On the branches they are in
\texttt{research/\allowbreak{}frontier/\allowbreak{}beyond/\allowbreak{}BEYOND\_\allowbreak{}TABLE.md} (each branch has its own version),
\texttt{research/\allowbreak{}frontier/\allowbreak{}beyond/\allowbreak{}BEYOND\_\allowbreak{}SMALL.md} (beyond-small), \texttt{research/\allowbreak{}frontier/\allowbreak{}beyond13/\allowbreak{}} (beyond-d13)
and \texttt{research/\allowbreak{}frontier/\allowbreak{}beyond11/\allowbreak{}} (beyond-d11). Master has no BEYOND\_TABLE.md yet. Where the branches and
master disagree, the table takes the largest value that has a check of its own, and says where it is. The
disagreements: (17,5), master 35,618, beyond-small 35,728; (3,11), master 1,344 (a voltage lift from own-small), every
BEYOND\_TABLE.md 1,474; (3,12), master 2,368, the BEYOND\_TABLE.md files 2,176; (17,7), master 5,693,959, the
BEYOND\_TABLE.md files 5,631,276; (12,12) and (13,12), master and beyond-d11b about 30 billion, the older tables about
20 billion; the D = 11, 12 values at Δ = 4..10, which the older tables give as found on October 26 and master holds at
the later, larger values; and (4,11), (5,11), (6,11), (4,12), (5,12), which beyond-d13 and beyond-d11b raised in
parallel, where master holds the larger value of each pair.

\hypertarget{the-table}{%
\section{3. The table}\label{the-table}}

\textbf{Columns.} \emph{construction}: ``Haar, orbit S'' is a Haar graph Haar(G, S) with S a union of orbits of a
group automorphism (section 5.1); ``Haar, random S'' has S searched with no imposed symmetry (section 5.2); ``(size
cap)'' marks the four cells where the job stopped at its size limit with the first S that worked. \emph{Group or base
graph}: Z\textsubscript{m} ⋊\textsubscript{r} Z\textsubscript{k} is the group of pairs (a, j), a mod m, j mod k, with
(a, j)(b, l) = (a + r\textsuperscript{j} b, j + l); ``Frobenius'' means m is prime and r has order k mod m (so the
group has trivial centre). \emph{Check}: ``+'' joins two independent checkers in one report (the harvest's or the
coordinator's); ``,'' joins checks run by the search session; a semicolon separates checks run at different times by
different people:

\begin{itemize}
\tightlist
\item
  gen = \texttt{research/\allowbreak{}frontier/\allowbreak{}haarcheck\_\allowbreak{}gen.c}; sl = \texttt{research/\allowbreak{}frontier/\allowbreak{}haarcheck\_\allowbreak{}sl.c}; psl =
  \texttt{research/\allowbreak{}structured/\allowbreak{}tools/\allowbreak{}haarcheck\_\allowbreak{}psl.c} (the reference checkers); h2 =
  \texttt{research/\allowbreak{}frontier/\allowbreak{}haarcheck2.c}; ind = \texttt{research/\allowbreak{}frontier/\allowbreak{}check\_\allowbreak{}haar\_\allowbreak{}indep.py}; bh2 =
  \texttt{research/\allowbreak{}frontier/\allowbreak{}build\_\allowbreak{}haar2.py}; mc = \texttt{research/\allowbreak{}frontier/\allowbreak{}beyond13/\allowbreak{}metacheck.py} or
  \texttt{metacheck2.py} (beyond-d13 branch). All of these build the group from its definition and run a BFS from
  (e,A) and from (e,B); see section 4.
\item
  apsp = all-pairs BFS with no symmetry assumed (\texttt{research/\allowbreak{}lifts/\allowbreak{}coordinator\_\allowbreak{}check/\allowbreak{}apsp.c}, or its
  multithreaded versions apsp\_mt / apsp\_fast); ``coord.'' = run by the coordinator on the Mac mini; ``8 ranges'' =
  the source vertices split over eight cloud workers.
\item
  auto = exact BFS from one vertex per orbit of an automorphism group that the checker verifies itself
  (\texttt{research/\allowbreak{}frontier/\allowbreak{}beyond/\allowbreak{}hex6/\allowbreak{}apsp\_\allowbreak{}auto.c}, beyond-small branch).
\item
  samp = BFS from 64 random vertices only (not a proof); ``theorem'' = the diameter bound is a theorem for every
  choice in the construction.
\item
  † = run by this job; logs in \texttt{research/\allowbreak{}paper/\allowbreak{}beyond\_\allowbreak{}checks/\allowbreak{}}.
\end{itemize}

\emph{Where}: ``master'' = harvested or recorded by the coordinator on master; ``branch'' = only on the job's branch.
Footnotes: ᵃ master has 35,618 (coordinator apsp, October 31). ᵇ no exact check (section 4). ᶜ master has 1,344. ᵈ one
checker so far (section 4). ᵉ a published census graph, CSS(6272,7), is larger (section 1). ᶠ the session's apsp logs
are on the beyond-d11b branch (\texttt{research/\allowbreak{}frontier/\allowbreak{}runs/\allowbreak{}cloudcheck/\allowbreak{}d4D11\_\allowbreak{}meta\_\allowbreak{}523\_\allowbreak{}29\_\allowbreak{}9\_\allowbreak{}b2h01.check.txt},
\texttt{d5D11\_\allowbreak{}meta\_\allowbreak{}5023\_\allowbreak{}31\_\allowbreak{}2886\_\allowbreak{}b5k05.check.txt}); master's harvest has bh2 only.

\begin{longtable}[]{@{}
  >{\raggedright\arraybackslash}p{(\columnwidth - 10\tabcolsep) * \real{0.0650}}
  >{\raggedleft\arraybackslash}p{(\columnwidth - 10\tabcolsep) * \real{0.1220}}
  >{\raggedright\arraybackslash}p{(\columnwidth - 10\tabcolsep) * \real{0.2276}}
  >{\raggedright\arraybackslash}p{(\columnwidth - 10\tabcolsep) * \real{0.2927}}
  >{\raggedright\arraybackslash}p{(\columnwidth - 10\tabcolsep) * \real{0.1951}}
  >{\raggedright\arraybackslash}p{(\columnwidth - 10\tabcolsep) * \real{0.0976}}@{}}
\toprule\noalign{}
\begin{minipage}[b]{\linewidth}\raggedright
(Δ, D)
\end{minipage} & \begin{minipage}[b]{\linewidth}\raggedleft
value
\end{minipage} & \begin{minipage}[b]{\linewidth}\raggedright
construction
\end{minipage} & \begin{minipage}[b]{\linewidth}\raggedright
group or base graph
\end{minipage} & \begin{minipage}[b]{\linewidth}\raggedright
check
\end{minipage} & \begin{minipage}[b]{\linewidth}\raggedright
where
\end{minipage} \\
\midrule\noalign{}
\endhead
\bottomrule\noalign{}
\endlastfoot
(17, 5) & 35,728 & route-chart lift + 1,936 hubs (88 Z₂₂-orbits) & (16,5) graph: K₃,₃ + doubled matching, Z₂₂, fibre
F₄⁴ & apsp, bfs\_all; apsp† & branch ᵃ \\
(18, 5) & 47,600 & multi-doubled-matching lift & W(13); 2 doubled matchings; Z₂ × Z₅ & apsp; coord. apsp & master \\
(19, 5) & 52,428 & doubled-matching lift & W(16); Z₆ & apsp; coord. apsp & master \\
(20, 5) & 75,000 & route-chart lift over F₅ & K₃,₃ + doubled matching, Z₂₀, fibre F₅⁴ & apsp; coord. apsp & master \\
(17, 6) & 2,236,962 & Moore graph & H(16) & classical; auto & branch \\
(18, 6) & 3,017,196 & Moore graph & H(17) & classical; auto & branch \\
(19, 6) & 3,203,968 & Theorem D, q = 17, m = 2 & H(17) doubled with levels and whiskers & theorem (Lean); auto &
branch \\
(20, 6) & 5,227,320 & Moore graph & H(19) & classical; auto & branch \\
(17, 7) & 5,693,959 & cover of H(13) with freed incidences + hubs & 7 copies of H(13); 32,356 + 30,327 hubs & coord.
apsp; apsp (8 ranges) & master \\
(18, 7) & 7,240,212 & level graph & 9 copies of H(13), 4 matchings, Z₉ & theorem; samp & branch ᵇ \\
(19, 7) & 11,184,810 & level graph & 5 copies of H(16), 2 matchings, Z₅ & theorem; samp & branch ᵇ \\
(20, 7) & 15,658,734 & level graph & 7 copies of H(16), 3 matchings, Z₇ & theorem; samp & branch ᵇ \\
(17, 8) & 41,373,072 & Haar, orbit S & SL(2,53) × Z₁₃₉ & sl + h2 & master \\
(18, 8) & 57,637,680 & Haar, orbit S & SL(2,61) × Z₁₂₇ & sl + h2 & master \\
(19, 8) & 86,317,296 & Haar, orbit S & PSL(2,103) × Z₇₉ & psl + ind & master \\
(20, 8) & 122,739,120 & Haar, orbit S & PSL(2,71) × Z₃₄₃ & psl + ind & master \\
(17, 9) & 501,707,400 & Haar, orbit S & PSL(2,101) × Z₄₈₇ & psl + ind & master \\
(18, 9) & 790,626,816 & Haar, orbit S & SL(2,127) × Z₁₉₃ & sl + h2 & master \\
(19, 9) & 1,218,394,320 & Haar, orbit S & SL(2,131) × Z₂₇₁ & sl + h2 & master \\
(20, 9) & 1,823,175,864 & Haar, orbit S & PSL(2,163) × Z₄₂₁ & psl + gen & master \\
(17, 10) & 7,022,052,336 & Haar, orbit S & PSL(2,313) × Z₂₂₉ & psl + gen & master \\
(18, 10) & 11,787,972,672 & Haar, orbit S & PSL(2,223) × Z₁₀₆₃ & psl + gen & master \\
(19, 10) & 19,233,722,496 & Haar, orbit S & SL(2,607) × Z₄₃ & sl + gen & master \\
(20, 10) & 30,518,778,360 & Haar, orbit S & PSL(2,421) × Z₄₀₉ & psl (2 halves) + gen & master \\
(3, 11) & 1,474 & Haar, all pairs \{x, y\} searched & Z₆₇ ⋊₉ Z₁₁ (Frobenius) & apsp; apsp† & branch ᶜ \\
(4, 11) & 30,334 & Haar, random S & Z₅₂₃ ⋊₉ Z₂₉ (Frobenius) & bh2, apsp ᶠ; mc† & master \\
(5, 11) & 311,426 & Haar, random S & Z₅₀₂₃ ⋊₂₈₈₆ Z₃₁ (Frobenius) & bh2, apsp ᶠ; mc† & master \\
(6, 11) & 2,122,518 & Haar, random S & Z₂₀₈₀₉ ⋊₁₂₉₇₂ Z₅₁ (Frobenius) & bh2; mc† & master \\
(7, 11) & 11,012,976 & Haar, orbit S & SL(2,53) × Z₃₇ & gen + sl & master \\
(8, 11) & 45,329,808 & Haar, orbit S & SL(2,283) & gen + sl & master \\
(9, 11) & 155,250,816 & Haar, orbit S & SL(2,223) × Z₇ & gen + sl & master \\
(10, 11) & 460,691,568 & Haar, orbit S & SL(2,613) & gen + sl & master \\
(11, 11) & 1,241,299,248 & Haar, orbit S & SL(2,853) & gen + sl & master \\
(12, 11) & 2,976,066,240 & Haar, orbit S & SL(2,239) × Z₁₀₉ & gen + sl & master \\
(13, 11) & 6,781,213,392 & Haar, orbit S & SL(2,563) × Z₁₉ & gen + sl & master \\
(14, 11) & 14,381,398,080 & Haar, orbit S & SL(2,1009) × Z₇ & gen + sl & master \\
(15, 11) & 28,655,908,512 & Haar, orbit S & SL(2,617) × Z₆₁ & gen + sl (2 halves) & master \\
(16, 11) & 29,961,590,400 & Haar, orbit S (size cap) & PSL(2,449) × Z₃₃₁ & gen only & branch ᵈ \\
(3, 12) & 2,368 & voltage lift & cubic 4 + 4-vertex base; Z₃₇ ⋊₆ Z₈ = SmallGroup(296,3) & coord. rebuild + apsp &
master \\
(4, 12) & 74,066 & Haar, random S & Z₁₂₇₇ ⋊₉₆₄ Z₂₉ (Frobenius) & apsp, mc; coord. apsp & master \\
(5, 12) & 1,058,342 & Haar, random S & Z₈₉₆₉ ⋊₂₀₅ Z₅₉ (Frobenius) & bh2; mc† & master \\
(6, 12) & 9,366,240 & Haar, orbit S & PSL(2,79) × Z₁₉ & gen + psl & master \\
(7, 12) & 59,838,504 & Haar, orbit S & PSL(2,67) × Z₁₉₉ & gen + psl & master \\
(8, 12) & 284,095,056 & Haar, orbit S & PSL(2,167) × Z₆₁ & gen + psl & master \\
(9, 12) & 1,105,785,120 & Haar, orbit S & PSL(2,241) × Z₇₉ & gen + psl & master \\
(10, 12) & 3,712,424,352 & Haar, orbit S & PSL(2,337) × Z₉₇ & gen + psl & master \\
(11, 12) & 11,189,148,096 & Haar, orbit S & PSL(2,223) × Z₁₀₀₉ & gen + psl & master \\
(12, 12) & 29,996,171,040 & Haar, orbit S & PSL(2,241) × Z₂₁₄₃ & gen + psl (2 halves) & master \\
(13, 12) & 29,997,644,040 & Haar, orbit S (size cap) & PSL(2,181) × Z₅₀₅₉ & gen + psl (2 halves) & master \\
(14, 12) & 29,996,171,040 & Haar, orbit S (size cap) & PSL(2,241) × Z₂₁₄₃ & gen + psl (2 halves) & master \\
(3, 13) & 3,400 & cyclic voltage lift & Desargues graph GP(10,3); Z₁₇₀ & apsp; coord. apsp & master \\
(4, 13) & 192,382 & Haar, random S & Z₂₂₃₇ ⋊₁₉₈₄ Z₄₃ (Frobenius) & apsp, mc; coord. apsp & master \\
(5, 13) & 3,594,910 & Haar, random S & (Z₁₁ ⋊₄ Z₅) × (Z₂₉₇₁ ⋊₁₇₆₄ Z₁₁), both Frobenius & mc; bh2† & branch \\
(6, 13) & 41,734,392 & Haar, orbit S & PSL(2,83) × Z₇₃ & gen; psl† & branch \\
(7, 13) & 312,227,760 & Haar, orbit S & SL(2,229) × Z₁₃ & gen; sl† & branch \\
(8, 13) & 1,739,514,000 & Haar, orbit S & SL(2,499) × Z₇ & gen; sl† & branch \\
(9, 13) & 7,924,036,032 & Haar, orbit S & SL(2,593) × Z₁₉ & gen; sl† & branch \\
(10, 13) & 19,980,294,768 & Haar, orbit S (size cap) & PSL(2,503) × Z₁₅₇ & gen only & branch ᵈ \\
(3, 14) & 5,706 & cyclic voltage lift & Pappus graph; Z₃₁₇ & apsp; coord. apsp & master ᵉ \\
(4, 14) & 497,042 & Haar, random S & Z₂₇₃₁ ⋊₂₄₀₀ Z₉₁ (Frobenius) & apsp, mc; coord. apsp & master \\
(5, 14) & 12,795,166 & Haar, random S & Z₁₄₈₇₈₁ ⋊₈₅₁₇₅ Z₄₃ (Frobenius) & mc; bh2† & branch \\
(6, 14) & 187,990,656 & Haar, orbit S & SL(2,97) × Z₁₀₃ & gen; sl† & branch \\
(7, 14) & 1,674,490,944 & Haar, orbit S & PSL(2,223) × Z₁₅₁ & gen; psl† & branch \\
(8, 14) & 11,325,238,056 & Haar, orbit S & PSL(2,397) × Z₁₈₁ & gen; psl† & branch \\
\end{longtable}

Counts: 44 of the 64 values are on master, each with two independent checkers or an all-pairs BFS: 24 at D = 11, 12,
the 12 Haar graphs at Δ = 17..20 and D = 8..10, 4 at D = 13, 14 with Δ = 3, 4, and (18,5), (19,5), (20,5), (17,7). For
four of them, (4,11), (5,11), (6,11) and (5,12), the second checker on master's file was run by this job. 20 are on
branches only: 10 at D = 13, 14 with Δ = 5..10 (nine now with two checkers, the second one run by this job; (10,13)
with one), (16,11) (one checker), the four D = 6 values (exact check via automorphisms), the three level graphs at D =
7 (theorem plus sampled BFS), and two values above master's, (17,5) and (3,11), each checked all-pairs twice. One
value, (3,14), is below a published census graph.

\hypertarget{how-the-values-were-checked}{%
\section{4. How the values were checked}\label{how-the-values-were-checked}}

\textbf{Haar graphs.} Haar(G, S) has vertices (g, A) and (g, B) for g in G and edges (g, A) \textasciitilde{} (sg, B)
for s in S. Right multiplication by any h in G maps edges to edges and acts transitively on each side, so every vertex
of one side has the same eccentricity, and the diameter is the larger of the eccentricities of (e, A) and (e, B). Two
BFS therefore decide the diameter exactly. Both are needed: for odd D and non-abelian G one side can pass while the
other fails (the first odd-D version of the new-cells-beyond search counted one side only; the checkers, which run
both BFS, rejected a (10,11) candidate it produced, and the file was discarded). Every checker listed above builds the
group from its definition (matrices mod p, or pairs (a, j) with the semidirect product rule), checks that the Δ
generators are distinct, and runs both BFS. The reference checkers sl and psl were written separately from gen, with
their own element indexing; above about 20 billion vertices their two BFS are run as two jobs (``2 halves''), which is
the same check split in two. bh2 can also write the edge list, so small Haar graphs can be given to apsp as well.

\textbf{Edge lists.} The non-Haar graphs (lifts, hub graphs, the level graphs at small q) were checked by apsp, which
reads an edge list and runs a BFS from every vertex (64 sources at a time) with no symmetry assumed; it also checks
simplicity, the maximum degree and bipartiteness. By the lab's convention each graph is also run at diameter D − 1,
where it must fail.

\textbf{D = 6 (Moore graphs and Theorem D).} These have 2.2--5.2 million vertices; an all-pairs BFS of each takes
hours on a 4-core cloud machine (the 5.7-million-vertex (17,7) graph was checked all-pairs that way, split over eight
machines). The beyond-small job checked them exactly with apsp\_auto: it is given permutations of the vertex set,
verifies that each maps every edge to an edge, and runs a BFS from one vertex per orbit of the group they generate
(eccentricity is constant on orbits). For H(q) the permutations come from diagonal matrices that preserve the quadric
and Tits' line equations (10,026 to 16,230 orbit representatives); for Theorem D at q = 17 one generator of order 16,
with 200,486 orbit representatives. apsp\_auto agrees with apsp on H(4) and on the Theorem D graphs at q = 5 and 7,
and fails at D − 1. The H(q) values are also classical (Tits), and Theorem D is proved in Lean for every odd q.

\textbf{D = 7 level graphs.} (18,7), (19,7), (20,7) are 2r + 1 copies of a generalized hexagon joined by r matchings.
Their diameter is at most 7 for every choice of matchings (research/synthetic/SYNTHETIC.md, Theorem 4.2), and the
builder was checked all-pairs at q = 3, 4 for r = 1..3. The instances themselves (7.2--15.7 million vertices) only had
a BFS from 64 random vertices, which found eccentricity 7. That is a sanity check of the builder, not a check of the
graph; the value rests on the theorem. No exact check (all-pairs or by automorphisms) has been run for these three
graphs.

\textbf{Second checks run by this job.} Where a row had only the search session's checker, this job ran a different
one on the same file (sha256 in each log): sl or psl on (6,13), (7,13), (8,13), (9,13), (6,14), (7,14), (8,14); bh2 on
(5,13) and (5,14); mc on master's (4,11), (5,11), (5,12) and (6,11), whose master logs have bh2 only (by the session
and again by the harvest); and apsp on the 1,474-vertex (3,11) graph and the 35,728-vertex (17,5) graph. All of them
reported VALID. The reference check of (10,13) (19,980,294,768 vertices) needs about 7 GB per half: the two halves did
not fit in this session's memory side by side (the A half was killed after 20 minutes when the B half started), and
running them one after the other would have taken about two more hours, so it was stopped unfinished. The reference
check of (16,11) (29,961,590,400 vertices, about 11 GB per half) was not attempted.

\hypertarget{methods}{%
\section{5. Methods}\label{methods}}

\hypertarget{haar-graphs-with-orbit-closed-generator-sets-ux3b4-17..20-at-d-8..10-ux3b4-6-or-7-at-d-11..14}{%
\subsection{5.1 Haar graphs with orbit-closed generator sets (Δ = 17..20 at D = 8..10; Δ ≥ 6 or 7 at D =
11..14)}\label{haar-graphs-with-orbit-closed-generator-sets-ux3b4-17..20-at-d-8..10-ux3b4-6-or-7-at-d-11..14}}

\textbf{Why Haar graphs, and how big.} For a Haar graph the diameter condition is a covering condition in G: every
element must be reached by a word of length at most D − 1 in S and S\textsuperscript{−1} (alternating). About w = Δ(Δ
− 1)\textsuperscript{D−2} such words exist. If they behaved like random elements, the expected number of elements
missed would be μ = N e\textsuperscript{−w/N} with N = \textbar G\textbar{} (research/frontier/SIZE\_LAW.md).
Annealing closes the last holes only when μ is small (SIZE\_LAW.md measured a limit μ* of about 9--40, depending on
effort), so the jobs chose group orders by this ``size law''. The Δ = 17..20 values were taken at μ ≈ 3--13 and not
pushed further; the D = 11..14 values with orbit S sit at μ ≈ 5--14 (the branch tables give μ for each row); the four
cells at the size cap have μ ≈ 0.

\textbf{Groups.} G = SL(2, p) × Z\textsubscript{n} or PSL(2, p) × Z\textsubscript{n}, p prime. These have many orders
near any target, cheap arithmetic, and an automorphism α of order 3: conjugation by a matrix of order 3 on the first
factor, times z ↦ rz with r of order 3 mod n on the second (for n = 1, or n with no such r, α acts on the first factor
only).

\textbf{Orbit search.} S = \{e\} ∪ F ∪ (orbits of α of size 3), with F a few α-fixed elements. Making S α-invariant
shrinks the search space and makes the whole word set α-invariant; in the lab's earlier runs at D = 7..10 this broke
the wall where random S stops working (README, research/structured/, research/theory/). The search anneals the orbit
representatives; the cost is the number of uncovered elements, counted by a bit-parallel BFS
(\texttt{research/\allowbreak{}bighaar/\allowbreak{}haar\_\allowbreak{}big.c} on the new-cells-beyond branch, and the beyond-d11 version
\texttt{research/\allowbreak{}frontier/\allowbreak{}beyond11/\allowbreak{}haar\_\allowbreak{}big0.c}, 64-bit indices, N/2 bytes of memory, so
10\textsuperscript{10}-element groups fit in about 5 GB). beyond-d11 added two things: S made of free orbits only,
without e (right for 3 \textbar{} Δ), and a fix in the choice of the automorphism for orbit sizes 4, 6 and 10 (the old
code could pick a matrix of order 2 or 3 there; orbits of size 3 and 5 were not affected, so no earlier result
changed).

\textbf{Odd D is harder than even D.} For D = 2h + 1 the diameter condition is
(S\textsuperscript{−1}S)\textsuperscript{h} = G and (SS\textsuperscript{−1})\textsuperscript{h} = G. The beyond-d11
job proved (research/frontier/beyond11/ODD\_D\_HOLES.md, on its branch) that the words reaching an involution come in
pairs (reverse and invert), and the words reaching a central element are closed under rotation by two letters, and a
central hole is a hole on both sides at once. In the random model an involution is then missed with probability about
e\textsuperscript{−λ/2} and a central element about e\textsuperscript{−λ/h}, against e\textsuperscript{−λ} for an
ordinary element (λ = w/N). PSL(2, p) has about p²/2 involutions, SL(2, p) only −I, and G × Z\textsubscript{n} has n
or 2n central elements. The rule that came out of this, and that gave most D = 11 and 13 values: at odd D use SL(2, p)
× Z\textsubscript{n} with small n; at even D, PSL(2, p) × Z\textsubscript{n} with moderate or large n is fine. The
exceptions in the table: (12,11), SL(2,239) × Z\textsubscript{109}, found before the rule; (16,11), PSL(2,449) ×
Z\textsubscript{331}, and (10,13), PSL(2,503) × Z\textsubscript{157}, both at the size cap, with λ = 41 and 32 (for
(16,11) beyond-d11b notes that λ is then large enough for the special elements to be reached too); and (6,13),
PSL(2,83) × Z\textsubscript{73} (beyond-d13 found that PSL with moderate n, 40 to 75, also worked at D = 13). PSL with
n = 1 is the worst choice at odd D: in the measured random S almost all holes (93--98 \%) were involutions, although
the same groups work at even D. Measured hole counts of random S were higher than these rates predict, by 1.3--1.5
times at μ ≈ 10 and 2--2.3 times at μ ≈ 30 (beyond-d11 card r3/e14).

\textbf{Size caps.} The jobs set a size limit so that the low-memory checkers finish in hours: 20 billion vertices for
beyond-d11 and beyond-d13, 30 billion for beyond-d11b. In (16,11), (13,12), (14,12) and (10,13) the first S tried
worked at the cap, so these values are limits of checking time, not of the search. (12,12) and (15,11) are also close
to the cap, but there the search had to work: μ ≈ 12 and 45 evaluations, μ ≈ 5 and 13 evaluations.

\hypertarget{metacyclic-haar-graphs-ux3b4-3..6-at-d-11..14}{%
\subsection{5.2 Metacyclic Haar graphs (Δ = 3..6 at D = 11..14)}\label{metacyclic-haar-graphs-ux3b4-3..6-at-d-11..14}}

At small Δ the groups are small (10\textsuperscript{3} to 10\textsuperscript{7} elements) and other groups can be
tried. Two kinds won.

\textbf{Odd-order and Frobenius groups.} In Z\textsubscript{m} ⋊\textsubscript{r} Z\textsubscript{k} with m prime and
r of order k \textgreater{} 1 mod m, the centre is trivial, and if mk is odd there are no involutions (proved by both
beyond-d13, \texttt{research/\allowbreak{}frontier/\allowbreak{}beyond13/\allowbreak{}THEORY.md} Fact 1, and beyond-d11b,
\texttt{research/\allowbreak{}frontier/\allowbreak{}beyond11/\allowbreak{}FROBENIUS.md} on the beyond-d11b branch, and its card r2/h17). So the special holes
of section 5.1 do not exist. These groups beat the SL/PSL values at odd D: (4,13) 192,382 against 170,520 for the best
PSL(2,29) × Z\textsubscript{7} orbit graph, (5,13) 3,594,910 against 3,217,056, and (4,11) 30,334 against 24,360.
Products of two such groups also work ((5,13)). The two jobs measured different reasons. beyond-d13 points to the
missing special holes. beyond-d11b (card r2/h13, at Δ = 4) found that the balls of radius 8 in the Frobenius groups
reach 0.99--1.00 of the random-words model and those of small SL(2, p) only 0.89--0.92, because short relations in
small SL(2, p) make words collide, and that involution and central holes were not the cause in that measurement. Both
effects may matter; we do not know their shares.

\textbf{Random S instead of orbits.} At Δ = 4 the orbit form \{e, x, αx, α²x\} has one free element, and at Δ = 5 one
fixed generator plus one orbit; a random S has Δ − 1 free elements. beyond-d13 gives this as the reason the metacyclic
groups with random S also won the even-D cells (4,12), (4,14) and (5,12), by 1.6--8 \% over the orbit graphs. At
(5,14) the runner-up already had random S (SL(2,97) × Z\textsubscript{7}, 12,776,064), and the metacyclic graph won by
only 0.15 \%. The search is plain annealing of S (\texttt{research/\allowbreak{}frontier/\allowbreak{}haar\_\allowbreak{}gen3.c}, and the beyond-d11b
versions with save/resume and a targeted finisher, \texttt{research/\allowbreak{}squads/\allowbreak{}beyond11b/\allowbreak{}tools/\allowbreak{}}; beyond-d13's
bit-parallel \texttt{research/\allowbreak{}frontier/\allowbreak{}beyond13/\allowbreak{}meta/\allowbreak{}hmeta.c}), plus scans of every metacyclic group in an order range
with short runs on each.

\textbf{Not too abelian.} Z\textsubscript{m} ⋊\textsubscript{r} Z\textsubscript{k} has an abelian subgroup of index k,
and for small k words collide. Measured at (4,13) with random S (squad card r2/e11 of beyond-d13): the words of length
12 from (e, A) reach 1 \% of G for k = 1, 6 \% for k = 3, 73 \% for k = 5, 91 \% for k = 9 and 97--99 \% for k ≥ 11.
The Δ = 4, 5, 6 values in the table use single groups with k = 29 to 91, or (at (5,13)) a product of two groups with k
= 5 and k = 11.

\textbf{Cubic cells.} At Δ = 3 with S = \{e, x, y\} all pairs (x, y) can be searched in each group, in 0.1--2 s per
group. The new-cells-beyond job swept every metacyclic group of order 541--1,100 at D = 11 (641 groups, one r of each
order per m) and found Z\textsubscript{67} ⋊\textsubscript{9} Z\textsubscript{11} (1,474 vertices), the only
metacyclic hit in that range. At D = 12 the best metacyclic Haar graph (2,176, Z\textsubscript{68} ⋊\textsubscript{3}
Z\textsubscript{16}) is below the own-small job's voltage lift over Z\textsubscript{37} ⋊\textsubscript{6}
Z\textsubscript{8} (2,368), which the table lists.

The metacyclic idea reached Δ = 6 only at D = 11; at (6,12), (6,13) and Δ = 7 it stopped at 2 to 14 uncovered
elements, and SL/PSL orbit graphs stayed ahead.

\hypertarget{cyclic-voltage-lifts-ux3b4-3-at-d-13-14}{%
\subsection{5.3 Cyclic voltage lifts (Δ = 3 at D = 13, 14)}\label{cyclic-voltage-lifts-ux3b4-3-at-d-13-14}}

Take a small cubic bipartite base graph B and a voltage x\textsubscript{e} in Z\textsubscript{n} on each edge; the
lift has vertices (v, g) and edges (u, g) \textasciitilde{} (w, g + x\textsubscript{e}). Rotation g ↦ g + c is an
automorphism, so only BFS from the base vertices (with g = 0) are needed. The beyond-d13 job annealed the voltages
(\texttt{research/\allowbreak{}frontier/\allowbreak{}beyond13/\allowbreak{}lifts/\allowbreak{}lift.c}, from squad card r2/e16; for short runs the cooling has to be spread
over about 1,000 evaluations, not the default 20,000) and then ran an exhaustive search over single-voltage changes
with a tabu list (\texttt{lifts/\allowbreak{}scan\_\allowbreak{}h02.c}, from card r5/h02). Results: the Desargues graph GP(10,3) over
Z\textsubscript{170} at D = 13 (3,400 vertices) and the Pappus graph over Z\textsubscript{317} at D = 14 (5,706). Both
beat every Haar graph tried at those cells (2,880 for A\textsubscript{6} × Z\textsubscript{4} at D = 13, all pairs
searched; 5,040 for PSL(2,7) × Z\textsubscript{15} at D = 14).

\emph{An upper bound for this method} (THEORY.md, Fact 2). The lift of a walk from (s, 0) ends at (v, Σ
c\textsubscript{e} x\textsubscript{e}), where c is the walk's net traversal vector (signed edge counts). So if the
lift has diameter ≤ D, then n is at most the number of distinct net traversal vectors of walks from s to v of length ≤
D, for every pair s, v. For GP(10,3) this allows n ≤ 647 at D = 13 and n ≤ 1,197 at D = 14; the lifts found have n =
170 and (for GP(10,3) at D = 14) 282. At D = 11, 12 beyond-d13's cyclic lifts did not beat the lab's existing values,
and at Δ = 4, 5 lifts of K\textsubscript{4,4}, PG(2,3), GQ(3,3) and random bases stayed below the Haar values.

\hypertarget{route-chart-lift-plus-hubs-175}{%
\subsection{5.4 Route-chart lift plus hubs: (17,5)}\label{route-chart-lift-plus-hubs-175}}

(17,5) was the one cell at D ≤ 7 in this note where none of the lab's formula families applies: the
multi-doubled-matching lifts need Δ = q + 1 + (an even number), and route charts need Δ = (letters) × (field size); 17
is neither for a usable prime power q. The new-cells-beyond job left only the floor 33,792, the (16,5) graph itself.
The beyond-small job added vertices to that graph.

\textbf{Route-chart lifts} (\texttt{research/\allowbreak{}theorems/\allowbreak{}chartlift/\allowbreak{}chartlib.py}). Start from a small bipartite
\emph{controller} graph whose edges are properly coloured with ℓ colours (\emph{letters}); letter c carries an
invertible matrix A\textsubscript{c} and a vector b\textsubscript{c} over F\textsubscript{q}. Each controller vertex u
becomes a \emph{fibre} of vertices (u, x), x in F\textsubscript{q}\textsuperscript{s}, and a controller edge u -- v of
letter c becomes the q\textsuperscript{s+1} edges (u, x) \textasciitilde{} (v, A\textsubscript{c}x +
βb\textsubscript{c}), β in F\textsubscript{q}. The lift has degree ℓq. Along a sequence of letters (a \emph{pattern})
the reachable set in the fibre is an affine subspace; the pattern is \emph{singular} if that subspace is not the whole
fibre, and a good chart has few singular patterns.

\textbf{Base.} The (16,5) graph (\texttt{research/\allowbreak{}theorems/\allowbreak{}graphs/\allowbreak{}build\_\allowbreak{}d16D5\_\allowbreak{}n33792.py}): K\textsubscript{3,3} with
a doubled perfect matching, lifted by Z\textsubscript{22} (66 + 66 controller vertices), 4 letters, fibre
F\textsubscript{4}\textsuperscript{4} (q = 4, s = 4). It has 33,792 vertices, is 16-regular, has diameter 5, and every
vertex has one spare degree when Δ = 17.

\textbf{Proposition (hubs).} Let G be bipartite with sides A, B and diameter 5. Add new vertices (hubs): a hub on side
A is joined to a set T ⊆ B, a hub on side B to a set T ⊆ A, each old vertex in at most one T. If (1) every vertex of
the hub's own side is within distance 3 of T, and (2) any two hubs on the same side have neighbours x, x′ with
d\textsubscript{G}(x, x′) ≤ 2, then the new graph has diameter ≤ 5.

\emph{Proof.} Old pairs are still within 5. A hub h on side A reaches all of A within 4 by (1), and all of B within 5,
since distances between B-vertices in G are at most 4 (G is bipartite with diameter 5). Two A-hubs are within 4 by
(2). An A-hub reaches every A-vertex within 4, so it reaches every B-hub, which has a neighbour in A, within 5. ∎

The two sides are independent, so the best packings of each side combine. With \textbar T\textbar{} = 16 the maximum
degree is 17.

\textbf{Hubs from fibre lines, in orbits.} A set T of 4 affine lines \{(v, y + βd) : β ∈ F\textsubscript{4}\} (d a
direction in the fibre) in 4 controller vertices covers a whole side within distance 3 for many choices (squad card
r1/c03). Translation g ↦ g + 1 on the controller coordinate is an automorphism of the base, so a good T gives 22 hubs
at once (card r2/d06), provided its translates are disjoint and pairwise compatible. Call the orbit of an old vertex
under the translation its \emph{key}; a side has 16,896 / 22 = 768 keys, the translates of T are disjoint exactly when
its 16 vertices have distinct keys, and orbits must not share keys. So a side holds at most 48 orbits, and the method
is capped at 33,792 + 2 · 48 · 22 = 35,904 vertices.

\textbf{Packing.} Greedy growth of orbit bases saturated at 32 orbits on side A and 37 on side B. Destroy-and-repair
loops (drop 2--5 orbits, regrow for 60 s, keep if not worse), an annealer over orbit packings, and ``region'' destroys
(drop every orbit with a key in a randomly chosen structured region, then regrow; card r4/f07) brought side A to 43
orbits and side B to 45: 88 orbits, 1,936 hubs, 35,728 vertices, 25.6 \% of the Moore bound. Twelve more search cards
in a fifth round found no 44th or 46th orbit. Master has the earlier 35,618 (83 orbits), checked all-pairs by the
coordinator; 35,728 is checked all-pairs by the session, by a separate plain BFS from every vertex
(\texttt{bfs\_\allowbreak{}all}), and by this job.

\hypertarget{the-other-rows}{%
\subsection{5.5 The other rows}\label{the-other-rows}}

\begin{itemize}
\tightlist
\item
  \textbf{(18,5), (19,5): doubled-matching lifts of W(q)} (\texttt{research/\allowbreak{}frontier/\allowbreak{}beyond/\allowbreak{}MDM.md} on the branches,
  \texttt{research/\allowbreak{}d10D7/\allowbreak{}DM\_\allowbreak{}D5.md}). With k doubled matchings over Z\textsubscript{2} × Z\textsubscript{2k+1} the
  lift of the symplectic quadrangle W(q) has degree q + 1 + 2k and 4(2k + 1)(q + 1)(q² + 1) vertices, and diameter ≤ 5
  as soon as one labelling of the flags by parities splits the q + 1 shortest paths between every two opposite
  elements (elements at distance 4 in W(q)); the matchings can be anything. (18,5): q = 13, k = 2; (19,5): q = 16, k =
  1 (Z\textsubscript{6}).
\item
  \textbf{(20,5): route-chart lift over F\textsubscript{5}.} The (16,5) construction with fibre
  F\textsubscript{5}\textsuperscript{4} and 4 letters × 5 scalars, controller K\textsubscript{3,3} + doubled matching
  lifted by Z\textsubscript{20} (60 controller vertices per side).
\item
  \textbf{D = 6.} H(q) for q = 16, 17, 19 are Moore graphs; (19,6) is Theorem D of the lab's polished notes (the
  hexagon H(17) doubled with levels and an exact whisker set), proved for odd q and checked in Lean.
\item
  \textbf{D = 7.} 2r + 1 copies of H(q) joined by r matchings (section 4); for (17,7) the conics-2 job freed some
  incidences of the 7-copy cover of H(13) and used the freed slots for 32,356 + 30,327 hub vertices joined by a hub
  graph of degree at most 10 (SYNTHETIC.md §4.2), 5,693,959 against 5,631,276 for the plain cover. For (18,7), (19,7),
  (20,7) the plain covers are listed: the new-cells-beyond job tried a simpler hub surgery, did not establish its hub
  condition for q = 13 and 16, and estimated its gain at 0.1--0.4 \%; the conics-2 construction (which gained 1.1 \%
  at (17,7)) was run for (17,7) only.
\end{itemize}

\hypertarget{what-failed}{%
\section{6. What failed}\label{what-failed}}

\begin{itemize}
\tightlist
\item
  \textbf{(17,5) before the hubs.} A fifth translation vector on one route-chart letter (17 = 5 + 4 + 4 + 4) cut the
  far pairs at Z\textsubscript{23} from 106,000 to 47,000 but not to 0, and cannot repair a singular pattern (section
  5.4) that uses the extended letter once. Lifts of W(13) with 3 extra degrees (8 sheets): with a doubled and a plain
  matching, the director's argument suggests that the conditions (needed for almost all pairs) force the labelling to
  be a coboundary, which kills the needed parity condition unless the matching is a polarity (a duality of order 2);
  on W(8), which has one, the far pairs fell from 174,000 to 2,432 but not to 0. In any cover of W(16) two copies of
  one vertex are at distance more than 5 (a path between them projects to a closed walk with nonzero voltage, which
  contains a cycle of W(16), and W(16) has girth 8). Haar graphs: best 20,200 (AGL(1,101)). A projective-plane
  scaffold for the hubs left 12,182 far pairs. Hub packing itself sits on a plateau at 43 + 45 orbits.
\item
  \textbf{Larger μ at D = 11.} (12,11) SL(2,601) × Z\textsubscript{7} (μ 10.7) stopped at 2 holes after about 230
  evaluations and SL(2,193) × Z\textsubscript{211} (μ 10.3) at 2 with −I its only hole; (11,11) PSL(2,257) ×
  Z\textsubscript{73} (μ 10.1) stopped at 9, while SL(2,853) at the same μ solved after 892 evaluations. The later job
  beyond-d11b did not try to raise (11..14, 11) or (11,12): the gains would be 1--3 \% and its machine was busy with
  the 30-billion-vertex checks.
\item
  \textbf{Near misses that did not close} (1--3 uncovered elements): (6,11) SL(2,103), (7,11) SL(2,67) ×
  Z\textsubscript{19}, (9,11) SL(2,137) × Z\textsubscript{31}, (6,12) PSL(2,211), (9,12) PSL(2,167) ×
  Z\textsubscript{241}, (6,13) SL(2,59) × Z\textsubscript{103} (the two holes are a transpose pair), (6,14) PSL(2,83)
  × Z\textsubscript{331}, (7,13) PSL(2,131) × Z\textsubscript{139}; and (9,13) SL(2,677) × Z\textsubscript{13} stopped
  at 12.
\item
  \textbf{Other shapes for S.} Orbits of size 2, 4, 5, 6 at D = 11; transpose-closed S; −I as a generator;
  orbit-closed S inside metacyclic groups; a meet-in-the-middle finisher. None beat the order-3 orbits (in SL/PSL) or
  random S (in metacyclic groups). Near the limit the annealer's 1- and 2-hole states seem to be true local minima:
  for one 1-hole (4,12) state all 105,609 single-generator swaps were tried and none closed it.
\item
  \textbf{Other groups and bases.} Metacyclic groups at Δ ≥ 6 (except (6,11)); Frobenius groups at (7,11) and (8,11)
  (2 and 8 holes); non-abelian voltage groups on GP(10,3) (2 holes at 3,300 vertices); in beyond-d11b's first round
  Sz(8), F\textasciitilde p\textsuperscript{3} ⋊ C\textsubscript{k}, bipartite Cayley graphs of PGL(2,p),
  PSL(2,p\textsuperscript{e}) with a field automorphism, AGL(1,q) products, and voltage lifts of PG(2,3) (19,500 at
  (4,11)).
\item
  \textbf{D = 8 from hexagons.} The 2013 table's D = 8 values follow a pattern (2j + 1)(⌊(q − 1)/2⌋ + 3(j −
  1))·\textbar H(q)\textbar{} (\texttt{research/\allowbreak{}theorems/\allowbreak{}REGULARITIES.md}). Its j = 2 case is the order of Gómez's
  compound graphs ΩH\textsubscript{q}, which as defined have diameter at least 9 (research/synthetic/SYNTHETIC.md §6;
  a diameter-8 version with a labelling for each copy is known only for q = 2, 3); for other j no construction is
  known. Card r1/c10 of new-cells-beyond found a Z\textsubscript{20} lift of H(3) plus two matchings with diameter 8
  (14,560, the pattern's value at q = 3, j = 2) by annealing, but q = 4 stalled at 150 holes. If the pattern could be
  realized for large q, it would give about 1.6--1.9 times the Haar values at D = 8 (for example about 2.2 ·
  10\textsuperscript{8} at (20,8) against 1.2 · 10\textsuperscript{8}).
\item
  \textbf{D = 5 otherwise.} Route charts over F\textsubscript{3} (6 letters) and F\textsubscript{2} (9 letters)
  reached only 29,160 at (18,5) in 25-minute runs. Blind annealing of voltages on W(13) stayed in the millions of
  uncovered pairs where the structured doubled-matching lift solves in seconds.
\item
  \textbf{The machines.} Four cores and 15 GB per cloud session. A 20-billion-vertex check was killed once for memory
  when a squad card's checker and a search ran beside it; squad cards were then capped at 0.5 GB, which also capped
  what they could search at (7,14) and (8,13). Checking time grows with p as well as N (haarcheck\_gen's field tables
  leave the L2 cache near p ≈ 850--1,000): the 14.4-billion-vertex (14,11) check took about 4.4 CPU-hours.
\end{itemize}

\hypertarget{loose-ends}{%
\section{7. Loose ends}\label{loose-ends}}

\begin{enumerate}
\def\labelenumi{\arabic{enumi}.}
\tightlist
\item
  The reference check of (16,11) and (10,13): \texttt{haarcheck\_\allowbreak{}psl\ FILE\ A} and \texttt{haarcheck\_\allowbreak{}psl\ FILE\ B}
  (research/structured/tools/haarcheck\_psl.c) on \texttt{d16D11\_\allowbreak{}psl\_\allowbreak{}449\_\allowbreak{}331\_\allowbreak{}orb3-5-1\_\allowbreak{}s1.txt} (beyond-d11b
  branch) and \texttt{d10D13\_\allowbreak{}psl\_\allowbreak{}503\_\allowbreak{}157\_\allowbreak{}orb3-3-1\_\allowbreak{}s1.txt} (beyond-d13 branch), about 11 GB and 7 GB per half.
\item
  An exact check of the three level graphs at D = 7, or the hub step of (17,7) applied to them.
\item
  The coordinator's own check and recording of the other 17 branch-only values: (17,5) 35,728 and (3,11) 1,474, which
  are above master's values, the D = 13, 14 values at Δ = 5..10, (16,11), and the four D = 6 values.
\item
  An apsp check of the census graph CSS(6272,7) at (3,14) (and CAT(3360,4) at (3,13)), and a literature check for all
  64 cells.
\end{enumerate}

\hypertarget{files}{%
\section{8. Files}\label{files}}

\begin{itemize}
\tightlist
\item
  Table rows and PDF: \texttt{research/\allowbreak{}paper/\allowbreak{}make\_\allowbreak{}beyond.py} (re-derives every vertex count and percentage and
  rewrites the table in BEYOND.md; \texttt{-\/\allowbreak{}-pdf} also runs pandoc and tectonic). This note:
  \texttt{research/\allowbreak{}paper/\allowbreak{}BEYOND.md}, \texttt{BEYOND.tex}, \texttt{BEYOND.pdf}.
\item
  Second checks of this job: \texttt{research/\allowbreak{}paper/\allowbreak{}beyond\_\allowbreak{}checks/\allowbreak{}} (README there).
\item
  Haar solution files: \texttt{research/frontier/sols/d\textless{}Δ\textgreater{}D\textless{}D\textgreater{}\_*.txt}
  (master for D ≤ 12 except (16,11); the beyond-d13 branch for D = 13, 14 (and
  \texttt{research/\allowbreak{}frontier/\allowbreak{}beyond13/\allowbreak{}sols/\allowbreak{}} there for the metacyclic (4,12) 74,066, of which master has only the edge
  list); the beyond-d11b branch for (16,11)). Rebuild and check one with
  \texttt{python3\ research/\allowbreak{}frontier/\allowbreak{}build\_\allowbreak{}haar2.py\ FILE} (small) or
  \texttt{cc\ -O3\ -o\ sl\ research/\allowbreak{}frontier/\allowbreak{}haarcheck\_\allowbreak{}sl.c\ \&\&\ ./\allowbreak{}sl\ FILE}.
\item
  Edge lists: \texttt{research/\allowbreak{}frontier/\allowbreak{}beyond/\allowbreak{}graphs/\allowbreak{}} (master and branches),
  \texttt{research/\allowbreak{}frontier/\allowbreak{}beyond13/\allowbreak{}graphs/\allowbreak{}} (beyond-d13), \texttt{research/\allowbreak{}squads/\allowbreak{}beyond-small/\allowbreak{}graphs/\allowbreak{}}
  (beyond-small); the files are gzipped; check with
  \texttt{cc\ -O3\ -o\ apsp\ research/\allowbreak{}lifts/\allowbreak{}coordinator\_\allowbreak{}check/\allowbreak{}apsp.c\ \&\&\ gunzip\ -k\ G.edges.gz\ \&\&\ ./\allowbreak{}apsp\ G.edges\ Δ\ D}.
\item
  Squad logs: \texttt{research/\allowbreak{}squads/\allowbreak{}LOG.md} and \texttt{research/\allowbreak{}squads/\allowbreak{}beyond*/\allowbreak{}} on each branch.
\end{itemize}

\end{document}
